\scanpage{57}
\mspage{24}
de l'ens à $N(x)+1$ éléments, où $N(x)=\operatorname{card}(\Sigma_{\unsure{<|\leq}x})$.
Ou encore, si on regarde l'application-but
\[
\Sigma^{<}(d)\longrightarrow\unsure{\Sigma}=\Sigma^{<}(0)
\]
\[
(x_0<-<x_d)\longmapsto x_d
\]
la fibre de $x\in\Sigma$ s'identifie à l'ens \deleted{des parties}
$\unsure{\mathfrak P|\mathcal P}_d(\{1,.-,N(x)\})$.
La structure relative de $\Sigma^{\unsure{<|*}}(1)$ sur $\Sigma=\Sigma^{<}(0)$
% Isolated dark written sign directly below the bare middle Sigma, between
% this line and the next: two short forked left strokes joined to a compact
% loop or oval at the right. Identity and mathematical attachment unassigned.
\illegible{1}
est donc celle d'un ens tot. ordonné relatif, qui
permet de reconstituer le système des $\Sigma^{*}(d)$
\deleted{avec} $(d\in\mathbb N)$ avec les applications semi-simpliciales
correspondant aux morphismes cofinaux entre
les ens tot. ordonnés-types.
% Slanted marginal note beside the preceding three lines; its words are
% not an insertion after semi-simpliciales. The final circled character is unclear.
\margin{Tout ceci est donc à faire indépendt de toute condition topologique \unsure{\textcircled{a}}.}

\textcircled{c} $\forall\ x\in\Sigma$, les germes \inserted{$B_\alpha$} de sous-espaces \inserted{bordants} de $\Sigma$ en $x$
correspondants, par \unsure{\textcircled{a}}, aux $z_\alpha\in\Sigma^{\unsure{<|*}}(1)$ au dessus
de $x$ (par $\Sigma^{\unsure{<|*}}(1)\overset{\underline b}{\unsure{-|\longrightarrow}}\Sigma$) sont mutuellement transver-
saux (et à fortiori distincts).
Les germes intersections de certains des $B_\alpha$
correspondent donc aux pts des $\Sigma^{\unsure{<|*}}(d)$ au dessus
de $x$ (par $\Sigma^{\unsure{<|*}}(d)\overset{\underline b}{\relbar\joinrel\relbar}\Sigma$), comme
images du \deleted{\unsure{dit}} germe de $\Sigma^{\unsure{<|*}}(d)$ en ces pts.
% Separate cancelled small looped glyph immediately before Les; no letter assigned.
\deleted{\illegible{1}} Les espaces $\Sigma^{*}(d)$ $(d\in\mathbb N)$ et les
applications $\tfrac12$ simpliciales cofinales entre eux, \deleted{\unsure{satisfa}}
% In the preceding cancellation a small final loop follows the tall f;
% its reading as a remains doubtful, and the abandoned word is not completed.
ou encore l'espace relatif $\Sigma^{\unsure{<|*}}(1)$ sur $\Sigma^{<}(0)$ et la
structure d'ordre relative, est décrit entièrement
% Pale bottom-left foliation, apparently another hand; the top authorial number is 24.
\margin{56}
